\begin{table*}[t]
\centering
\small
\caption{Defense taxonomy in \asb{}. Mechanism types: \textbf{P}rompt-level, \textbf{T}ool-gating, \textbf{C}ontent-level, \textbf{M}ulti-signal. D10 (\civ{}) is our novel contribution.}
\label{tab:taxonomy}
\begin{tabular}{llllccc}
\toprule
ID & Name & Type & Key Mechanism & Modifies & Gates & Filters \\
   &      &      &               & Prompt   & Tools & Output \\
\midrule
D0  & Baseline           & ---  & No defense                          & --- & --- & --- \\
D1  & Spotlighting        & P    & Delimiter-based source marking      & \checkmark & & \\
D5  & Sandwich            & P    & Goal-reminder wrapping              & \checkmark & & \\
D7  & Input Classifier    & P    & Injection pattern removal           & \checkmark & & \\
\midrule
D2  & Policy Gate         & T    & Risk-level + whitelist enforcement  & & \checkmark & \\
D3  & Task Alignment      & T    & LLM goal--action consistency check  & & \checkmark & \\
D4  & Re-execution        & T    & Clean re-run comparison             & & \checkmark & \\
D8  & Semantic Firewall   & T    & Embedding-based drift detection     & & \checkmark & \\
D9  & Dual-LLM            & T    & Two-model screening                 & & \checkmark & \\
\midrule
D6  & Output Filter       & C    & Regex-based data leak detection     & & & \checkmark \\
\midrule
D10 & \textbf{CIV (ours)} & M    & Provenance + embed.\ compat.\ + plan dev. & & \checkmark & \\
\bottomrule
\end{tabular}
\end{table*}
